\section{Deeper confluent certificates and exact sparse families}\label{sec:confluent}

The incoming proof of the last quintic case uses scalar congruences and divided differences to make a finite test valid for arbitrarily large exponents. Its stated sufficient condition is $p>k+1$, where $k$ is the total number of powers. Retaining a factorial numerator gives the following stronger statement.

\begin{theorem}[Higher-depth, blockwise certificate]\label{thm:confluent}
Let $p$ be any prime, $s,m\ge1$, and let $A$ be a finite free commutative $\Z_{(p)}$-algebra. Suppose
\[
                    x^m=a+p^sD,\qquad p\nmid a,
\]
with $x,D\in A$ and $a\in\Z$. Let $\mathcal E$ be a finite set of distinct nonnegative integers, and let $n_r$ count its elements congruent to $r$ modulo $m$, for $0\le r<m$. Assume
\begin{equation}\label{eq:block-bound}
                     n_r<p^s\qquad\text{for every }r.
\end{equation}
If the columns
\[
                x^rD^j,\qquad 0\le r<m,\quad0\le j<n_r,
\]
are linearly independent in $A/pA$ over $\F_p$, then the powers $x^e$, $e\in\mathcal E$, are linearly independent in $A\otimes\Q$ over $\Q$.
\end{theorem}

\begin{proof}
Within a residue block write the exponents as $e_i=r+mt_i$. Divide the associated columns by $a^{t_i}$. They become
\[
                         x^r(1+p^sD/a)^{t_i}.
\]
For $j=0,\ldots,n_r-1$, take the Newton divided difference of order $j$ at the first $j+1$ distinct integers $t_i$, multiplied by $j!/p^{sj}$. These are invertible triangular column operations over $\Q$.

The coefficient of $(D/a)^i$ in the new column is
\[
 j!p^{s(i-j)}
       [t_0,\ldots,t_j]\binom{T}{i}.
\]
It is zero for $i<j$ and one for $i=j$. For $i>j$, its valuation is at least
\begin{equation}\label{eq:factorial-bound}
                       s(i-j)-v_p(i!/j!).
\end{equation}
Indeed, $i!\binom{T}{i}$ is an integer polynomial, and divided differences of integer polynomials at integer arguments are integers. For a monomial this is the complete homogeneous symmetric-polynomial identity, and linearity gives the general case.

Let $S_p(n)$ denote the base-$p$ digit sum. Since $j\le n_r-1\le p^s-2$,
\[
 S_p(j)\le s(p-1)-1,\qquad S_p(i)\ge1.
\]
Set $d=i-j\ge1$ and $C=s(p-1)-1\ge0$. Legendre's formula gives
\begin{align*}
 (p-1)\bigl(sd-v_p(i!/j!)\bigr)
 &=Cd-S_p(j)+S_p(i)\\
 &\ge C(d-1)+1>0.
\end{align*}
Every later coefficient is therefore divisible by $p$. The transformed column is integral and reduces to $x^r(D/a)^j$. Its factor $a^{-j}$ is a nonzero scalar modulo $p$.

A rational dependence between integral columns in a finite free $\Z_{(p)}$-module can be scaled to have integral coefficients with at least one unit. Reduction would contradict the assumed independence. Reverse the invertible column operations to obtain the result.
\end{proof}

The proof uses finite binomial expansions; no $p$-adic exponential or logarithm is required. Thus $p=2$ causes no convergence problem. The bound is on each residue block, not on total support. At depth one it needs $n_r<p$, already weaker than the preceding total-support condition.

\subsection{Sharpness of the strict threshold}

The strict inequality in \eqref{eq:block-bound} cannot be replaced universally by a weak inequality. Put $k=p^s$ and
\[
 A=\Z_{(p)}[D]/(D^k+D^{k-1}-1),\qquad x=1+kD,\quad m=1.
\]
The $k$ confluent columns $1,D,\ldots,D^{k-1}$ form a basis modulo $p$. Nevertheless
\[
 F(X)=(X-1)^k+k(X-1)^{k-1}-k^k
\]
vanishes at $x$ and has exactly $k$ nonzero coefficients. The coefficient of $X^{k-1}$ is zero; all the coefficients at exponents $0,1,\ldots,k-2,k$ are nonzero. At $p=3,s=1$, the example becomes the irreducible relation $x^3-3x-25=0$ after $D^3+D^2-1=0$.

This proves sharpness in the theorem's finite-free-algebra generality. It does not assert that the chosen order equals $\Z_{(p)}[x]$; in fact that narrower equality need not hold.

\subsection{Two exact seven-term families}

\begin{theorem}[Powered singular moduli]\label{thm:powered}
Let $\alpha$ be any singular modulus of discriminant $-87$, and let $\beta$ be any singular modulus of discriminant $-152$. Then, for every integer $t\ge1$,
\[
                     s(\alpha^{63t})=7,\qquad
                     s(\beta^{182t})=7.
\]
The exponents and rational coefficients of the annihilating polynomials are unrestricted.
\end{theorem}

\begin{proof}
Both ring class polynomials have degree six. Their exact integer coefficients and the certificates below are in \code{checks/confluent_certificates.json}. Their identification as class polynomials is inherited from the prior report's interval-certified CM calculations, whose complete verifier was rerun in this audit.\cite{prior_sparse}

In $A=\Z_{(2)}[X]/(H_{-87})$, exact reduction gives
\[
 x^{63}\equiv 1+16(x+x^2+x^3)\pmod {32A}.
\]
Thus $x^{63}=1+16D$ with $D\bmod2=x+x^2+x^3$. The six columns $1,D,\ldots,D^5$, represented with entries $0,1$ in the power basis, have integer determinant $-1$. They are independent modulo $2$, and $6<2^4$. Apply Theorem~\ref{thm:confluent} to any six distinct exponents divisible by $63$.

Likewise, in $\Z_{(3)}[X]/(H_{-152})$,
\[
 x^{182}\equiv26+9(2x^2+x^3+2x^5)\pmod {27A}.
\]
For $D=(x^{182}-26)/9$, the six columns $1,D,\ldots,D^5$ have the following representative matrix modulo $3$:
\[
 \begin{pmatrix}
 1&0&2&2&0&0\\
 0&0&2&0&2&2\\
 0&2&2&1&1&2\\
 0&1&0&0&0&0\\
 0&0&0&0&2&1\\
 0&2&1&2&0&1
 \end{pmatrix}.
\]
Its integer determinant is $-2$, hence is nonzero modulo $3$, and $6<3^2$.

For either powered number, a rational relation with at most six terms would select at most six such powers. A smaller dependent set could be enlarged to six distinct exponents, contradicting the established independence. On the other hand, the number lies in a six-dimensional $\Q$-algebra, so seven powers are dependent. This gives support exactly seven. Replacing $63$ or $182$ by any positive multiple retains the argument.
\end{proof}

These are exact infinite-exponent statements, not bounded searches. The modular powers were verified both by polynomial multiplication and by an independently written basis-shift recurrence. The matrices and determinants were checked with integer arithmetic.

The result does not prove $s(-87)=7$ or $s(-152)=7$. A six-term relation for the original singular modulus could involve several different exponent residue classes. That remaining possibility is precisely what a complete residue-cover search would have to exclude.

\subsection{Why this is a useful next step}

The local idea fits the source's arithmetic theme: retain divisibility which a coarser size estimate discards, then certify independence through a determinant. Here the implementation is elementary and the proofs are short enough to formalize independently. It gives a concrete partial result in the next support problem, while avoiding any claim that an analogous theta-lift or Selmer construction has already been built for a new $L$-function.

The deeper certificate should be included in a future search over six-term exponent patterns, especially at small primes. The output of that search must be an exhaustive finite cover of all residue patterns, including repeated residues; checking large exponent ranges alone would not prove the unbounded statement.
